\cotefile{68}
\scanpage{17}
\diagram{\begin{tikzcd}[ampersand replacement=\&,row sep=small,column sep=small]
 E_i=[-i,i] \& 0\leq i\leq m \& \text{NB. } E_m=G,\ E_1=\varepsilon\\
 F_j=[m-j,m+1\unsure{-|+}j] \& 0\leq j\leq m-1 \& F_m=E_m=G\\
 F_{\unsure{i|j}}*F_j=\displaystyle\sum_{0}^{m}\deleted{\unsure{\alpha}}\beta_{ij}^{k}E_{h} \& \beta_{ij}^{h}\geq0\\
 E_i*E_j=\displaystyle\sum\alpha_{ij}^{h}E_h\arrow[uu,bend left=63,start anchor=west,end anchor=south west] \& \alpha_{ij}^{h}\geq0
\end{tikzcd}}
\[
 E_i*F_j=\sum\gamma_{ij}^{k}F_h \qquad\qquad
 \unsure{\varepsilon|\zeta|s}=\unsure{\zeta|s|\xi|\varepsilon}^{-1}\quad\deleted{\unsure{0}}
\]
\[
 \mathcal A(\underline{g,\deleted{\unsure{g}}\varepsilon})=
 \begin{cases}
 g_2*\mathcal A(0)&\text{si }\varepsilon=+1\\
 g_2*\deleted{\unsure{\mathcal A}}\mathcal B(0)&\text{si }\varepsilon=-1
 \end{cases}
\]
\[
 \mathcal A(\tilde g)*\mathcal A(\tilde h)\unsure{\subseteq|\subset}\mathcal A(\tilde g\cdot\tilde h)
\]

\deleted{$\unsure{G'|\text{Soit}}\quad G'$ groupe \quad $G'\xrightarrow{\unsure{\lambda}}G$}

\margin{$\unsure{+|f}\quad A=g*\check A=\check A*g$\\
$B=h*\tilde B=\check B*h$\\
$A*B=g*\tilde A*h*\tilde B=g*\check A*\tilde B*h$}

$\tilde G$ gr\unsure{oupe}, $\tilde G\xrightarrow{\unsure{y|p|\nu|2|\eta}}G$ hom,
\inserted{$i:{}_2G\simeq\operatorname{Ker} y$}

\[
 \forall\tilde g\in G\qquad \mathcal A(\tilde g)\subset\mathcal P_f(G),
\]

\margin{On pose $C(\tilde g)$ \unsure{cône|c\=ome|comme} des $\displaystyle\sum_{A\in\mathcal A(\tilde g)}\alpha_A A\deleted{\unsure{*}},\ \alpha_A\geq0$}

a) $\mathcal A(\tilde g)$ tot ordonné

b) $A\in\mathcal A(\tilde g)\ \Longrightarrow\ A=\deleted{\unsure{g(\tilde g)}}g*\check A$ où $g=y(\tilde g)$
\margin{$\forall\tilde g\in\tilde G$}

c) $A\in\mathcal A(\tilde g),\ B\in\deleted{\unsure{\mathcal B}}\mathcal A(\tilde h)\ \Longrightarrow\ A*B\in C(\tilde g\tilde h)$
\margin{$\forall\tilde g,\tilde h\in\tilde G$}

\inserted{d) Si $u\in{}_2G,\ \tilde g\in\tilde G$, on a.}

\deleted{d) $\mathcal A(\tilde g\,u^2)=u*\mathcal A(\tilde g)$}\quad
$\deleted{\mathcal A(\unsure{u\tilde g*})=}\mathcal A(i(u)\tilde g)=u*\mathcal A(\tilde g)$

d) $\deleted{\unsure{\exists}}\forall\tilde g,\tilde h\in\tilde G\quad\exists s\in G$ tel que $\mathcal A(\tilde g)\cup\varepsilon_s\mathcal A(\tilde h)$ tot. ordonné

[ i.e. $\bigcup\unsure{\mathcal A k g|\mathcal A(\tilde g)|\mathcal A_{\tilde g}}$
% archive sheet number: 16
\scanpage{18}
% archive sheet number: 17
\scanpage{19}
\textcircled{1}\quad $|f,g|$

\deleted{$\unsure{\lvert f,g\rvert}$} $G$ un ensemble, $C(G)=\mathbb R[G]=\mathbb R^{(G)}$, (fonctions réelles à support fini), $\mathbb R[G]^+=\{f\in\mathbb R[G]\mid f\geq0\}$

\deleted{$f,g\in$} Éparpillement $f'$ de $f$ (ou $f,f'$ équivalentes, on écrit $f'\sim f$) sss $\exists\sigma\in \unsure{\mathfrak S|\mathcal S|S}_G,\quad f'=f\circ\sigma\ \Longleftrightarrow$ la suite des valeurs $\ne0$ de $f$, \deleted{\unsure{comptées|complètement}} rangées en décroissant et avec multiplicités, sont les $\overline{\mathrm m}$s $\Longleftrightarrow f(\mu)=f'(\mu)$ (image directe de la mesure canonique).

% slanted
\margin{Notation $A(f)$; \unsure{condition|aussi bien|ordonnée} $A(f')=A(f)$}

Si $f,g\in\mathbb R[G]^+$, on pose
\[
 |f,g|=\underset{f'\sim f}{\operatorname{Sup}}\langle f',g\rangle
 =\underset{g'\sim g}{\operatorname{Sup}}\langle f,g'\rangle
 =\underset{\substack{f'\sim f\\g'\sim g}}{\operatorname{Sup}}\langle f,g\rangle
\]
\margin{\uline{NB} $|f,g|=|f',g'|$ si $f'\sim f,\ g'\sim g$}

\deleted{\uline{\unsure{Proposi}}} Pour que $\langle f',g'\rangle$ réalise le Sup, il f. et s. que pour $\forall\alpha,\beta\in\mathbb R^{+*}$, $E_\alpha(f')$ et $E_\beta(g')$ soient en relation d'inclusion (i.e., $\subset$ ou $\supset$), i.e il f et s. que $E_*(f')\cup E_{\alpha}(g')$ soit tot ordonné $\Longleftrightarrow\exists$ un ordre total sur $G$ [Si $G$ est fini ou dénombrable] $\sim[0,n]$ ou $N$, tel que $f$ et $g$ soient décroiss\deleted{\unsure{antes}}

% slanted
\margin{\uline{NB} Si $f\sim g$ alors le Sup est réalisé sss $f'=g'$ et $|f,g|=|f,f|=|g,g|=\|f\|_2^2=\int f^2$}

\uline{NB} On pose $E_\alpha(f)=\{s\in G\mid f(s)\geq\alpha\}$ (\deleted{\unsure{donc}} $\alpha\leq\beta\Rightarrow E_\alpha(f)\supset E_\beta(f)$) et $E_*(f)=\{E_\alpha(f)$

\deleted{\uline{\unsure{Proposit}}} $f'\sim f\Longleftrightarrow\forall g\in C(G)^+,\ |f',g|=|f,g|\ \Longleftrightarrow\ \forall A\subset\mathcal P_f(G),\ |f',A|=|f,A|$

\textcircled{2}\quad \uline{$f'\prec f$}

\uline{Prop.} Conditions équivalentes sur $f',f\in\mathbb R[G]^+$

a) $\forall g\in\mathbb R[G]^+$, on a
\[
 |f',g|\leq|f,g|
\]

b) $\forall\deleted{\unsure{E}}A\in\mathcal P_f(G)$, on a
\[
 |f',\deleted{\unsure{E}}A|\leq|f,\deleted{\unsure{E}}A|
\]

c) $\forall\deleted{\unsure{\text{pour}}}\unsure{n|p}\in\mathbb N^{*}$, on a
\[
 \sum_0^{n} A(f')_i\leq\sum_0^{n} A(f)_i
\]
\margin{On écrit alors $f'\prec f$}
% archive sheet number: 18
\scanpage{20}
\uline{NB} C'est une relation \deleted{d'ordre}\inserted{de préordre}. La relation d'équiv. associée est $f'\sim f$

\margin{\uline{NB} Le Sup de $\langle f',g\rangle$ pour $f'\prec f$ est atteint \deleted{pour} \inserted{des} $f'\sim f\ldots$}

\deleted{\uline{Cor} $f'\prec f,\ g'\prec g\ \Rightarrow\ |f',g'|\leq\lvert f,g$}

\uline{Cor} $f'\prec f\ \Longleftrightarrow\ \forall g',g\in C(G)^+$, avec $g'\prec g$, on a
\[
 |f',g'|\leq|f,g|
\]

\uline{3}\quad \uline{\deleted{\unsure{Paires|P}}Couples prétassés}

À partir de maintenant, $G$ est un groupe, d'où opération $f*g$ sur $\mathbb R[G]^{*}$, $\mathbb R[G]^+$ est stable

\uline{Prop} $f,g\in\mathbb R(G)^+$, Conditions équivalentes

% margin note written sideways
\margin{\uline{NB} Soient $f,g\in\mathbb R[G]^+$. Considérons les $\deleted{\unsure{\Vert}}\|f'*g\|_{\infty}$ ou $f'*g(0)$ pour $f'\prec f$, alors le Sup est \deleted{\unsure{atte|atteint} \unsure{et vaut}}\inserted{$=$}, et il est atteint pour $f'\sim f$}

a) $f'\prec f,\ g'\prec g,\qquad \|f'*g'\|_\infty\leq\|f*g\|_\infty$

b) $\forall f'\prec f,\ g'\prec g\qquad (f'*g')(0)\leq\operatorname{Sup}_{s\in G}(f*g)(s)$
\[
 \begin{aligned}
 &\parallel &&\parallel\\
 &\langle f',\check g'\rangle &&\deleted{\unsure{\langle f,\check g\rangle}}\;\langle f\deleted{\unsure{*}},\varepsilon_s*\check g\rangle
 \end{aligned}
\]

a$'$) \deleted{$\unsure{g',f}\ \unsure{\forall f'}\ \unsure{\forall g'}$} a) b) avec $f'\sim f,\ g'\sim g$

c) $|f,g|=\|f*g\|_\infty$

d) $\exists s\in G$ tel que $E_\alpha(f)\cup E_*(\varepsilon_s\check g)$ tot. ordonné i.e.
\margin{$\Updownarrow\ E_\alpha(g)\cup E_\alpha(\check f*\varepsilon_s)\deleted{\unsure{E_*(\check g')}}$ tot. ordonné}
\[
 \forall\alpha,\beta\in\mathbb R^{+*},\qquad E_\alpha(f)\ \text{et}\ E_\beta(\varepsilon_s\check g)\ \text{en relation d'inclusion}
\]

\uline{Déf} On dit alors que le \deleted{\unsure{fon}}couple $(f,g)$ est \uline{prétassé}

( [Il n'est pas clair que la relation soit symétrique -- \deleted{cela signifie $\forall s$} mais OK si $G$ commutatif] )

\uline{Cor} $f\in\mathbb R[G]^+$ \deleted{symétrique} $(f,f)$ est prétassé $\Longleftrightarrow$ \deleted{\unsure{$f$ et $f$|$f$ et $\check f$}} \inserted{$\exists s\in G$} tel que $\check f=\varepsilon_s*f\quad(\Longleftrightarrow\check f=f*\varepsilon_s)$

\uline{NB} Si dans la prop. $g$ satisfait la condition précédente, alors la condition de la prop signifie aussi e) $\exists\unsure{t|s}\in G$ t.q. $E_\alpha(f)\cup E_\alpha(\varepsilon_t*g)$ \inserted{tot. ordonné}
\margin{$\Updownarrow\ E_\alpha(\varepsilon_{\unsure{t^{-1}|t|s}}*f)\cup E_\alpha(g)$ tot. ord.}
% archive sheet number: 19
\scanpage{21}
\deleted{\uline{\unsure{4.|5.}}} Donc si $f\unsure{{}_+|+}$ et $g$ sont sym. à translation près, la condition $(f,g)$ prétassé est une condition \uline{symétrique} en $f,g$

4. \uline{Couples tassés}.

$f,g\in\mathbb R[G]^+$.

\margin{Cette condition est inv. par translation. \deleted{\unsure{Donc}} \inserted{Si} $(f,g)$ est tassé, $\check g,\check f$ aussi. Donc si $f,g$ \deleted{\unsure{t}} sont sym. mod translation, alors $(f,g)$ tassé sss $(g,f)$ tassé}

\uline{Prop} Conditions équivalentes

a) $f'\prec f,\ g'\prec g\ \Longrightarrow\ f'*g'\prec f*g$

b) $f'\sim f,\ g'\sim g\ \Longrightarrow\ f'*g'\prec f*g$

c) $\forall h\in\mathbb R[G]^+$
\[
 \begin{aligned}
 \underset{\substack{f'\prec f\\g'\prec g\\h'\prec h}}{\operatorname{Sup}}(f'*g'*h')(0)
 &\leq\underset{\substack{\deleted{\unsure{h'\prec h}}\\h'\sim h}}{\operatorname{Sup}}f*g*h'(0)\\
 &\hphantom{{}\leq{}}\parallel\\
 &\left(\underset{h'\sim h}{\operatorname{Sup}}\|f*g*h'\|_\infty\right.
 \end{aligned}
\]

c$'$) $\forall\deleted{\unsure{E}}A\in\mathcal P_f(G)$,
\[
 \begin{aligned}
 \underset{\substack{A'\sim A\\f'\sim f\\g'\sim g}}{\operatorname{Sup}} f'*g'*A'(0)
 &\leq\underset{A'\sim A}{\operatorname{Sup}}f*g*A'(0)\\
 &\hphantom{{}\leq{}}\parallel\\
 &\left(\underset{A'\sim A}{\operatorname{Sup}}\|f*g*A'\|_\infty\right.
 \end{aligned}
\]

Cel implique que $(f,g)$ est \uline{prétassé}

\uline{Cor} Soit $S\subset\mathbb R[G]^+$, \uline{$S$ stable par $*$}. \deleted{Pour que}\inserted{Conditions équiv. :}

\deleted{\unsure{$S$ soit}} \inserted{a)} tt couple $(f,g)\in S\times S$ soit tassé, \deleted{il \unsure{f. et s.}}

b) \deleted{que.} Pour \deleted{tout} $n\in\mathbb N,\ n\geq2,\quad f_1,\text{--}\,f_n\in S,\quad f'_1,\text{--}\,f'_n\in\mathbb R[G]^+$ $f'_1\prec f_1$ --, $f'_n\prec f_n$ implique
\[
 f'_1*\text{--}*f'_n\prec f_1*\text{--}*f_n
\]

\margin{Condition K $\uparrow\downarrow$ \unsure{hypoth|hypoth.|hyp.}}
c) [ Si $\forall n\in\mathbb N^{*},\ \inserted{n\leq\operatorname{card}G,}\ \exists A\in S\cap\mathcal P_f(G)$ tel que $\operatorname{card} A=n$ ] $\longrightarrow$ ``Condition K''

$f'_1\prec f_1,\text{--},\ f'_n\prec f_n\quad(f'_i\in\mathbb R(G)^+,\ f_i\in S)\ \Longrightarrow$
\[
 \|f'_1*\text{--}*f'_n\|_\infty\leq\|f_1*\text{--}*f_n\|_\infty
\]
(ou $f'_1*\text{--}*f'_n(0)$

\deleted{Condition K\quad d) $f'_1\sim f_1,\ f'_2\sim f_2,\ \unsure{\text{--}}\ f'\sim f,\ g'\sim g,\ A'\sim A\in P_f(G)\cap S\ \Rightarrow\ f'*g'*A'(0)\leq\|f*g*A\|_\infty$}
% archive sheet number: 20
\scanpage{22}
% slanted
\margin{si (K) et si les couples $\in S\times S$ sont \uline{pré}tassés $\updownarrow$ \unsure{id|id.}}

d) $\forall A,B\in\mathcal P_f(G)\cap S,$ $A'\deleted{\unsure{\prec}}\sim A,\ B'\deleted{\prec}\sim B$, on a
\[
 A'*B'\prec A*B\qquad\text{i.e. \uline{$(A,B)$ tassé}}
\]
e) $\forall A,B,C\in\mathcal P_f(G)\cap S,\quad A'\sim A,\ B'\sim B,\ C'\sim C$, on a
\[
 \|A'*B'*C'\|_\infty\leq\|A*B*C\|_\infty
\]

\noindent\rule{\linewidth}{0.4pt}

\deleted{Supposons}

\unsure{Quitte|Quitt|Conclut} \unsure{à|:} saturer $S$, \unsure{ops} $S$ stable par translation

$S_0=S\cap\mathcal P_f(G)$ \unsure{ops} $S_0\not\ni\varnothing$

Soit $f\in S$. Alors
\[
 \begin{aligned}
 \text{a)}\quad&\exists t\in G,\quad f=\varepsilon_t*f
 &&[(f,f)\ \inserted{\text{pré}}\text{tassé}]\\
 \text{b)}\quad&\text{les } E_\alpha(f)\in S_0
 &&[(f,A)\ \inserted{\text{pré}}\text{tassé pour }A\in S_0]
 \end{aligned}
\]

\deleted{\unsure{Inv., soient} $f\in\mathbb R[G]^+$}

Soit $\deleted{\unsure{g|S}}\tilde S=\{f\in\mathbb R(G)^+\mid\text{satisfont a), b)}\}$

\deleted{\unsure{$f,g$|$f\in\tilde S$}} i.e
\[
 \left\{\begin{aligned}
 f&=\sum_I\alpha_i A_i&\alpha_i>0\quad A_i\in S_0\quad (A_i)\ \text{tot ord.}\\
 \deleted{\unsure{g}}\quad&\exists s\in G,& A_i=\deleted{\unsure{\check A_i}}\varepsilon_s*\check A_i\quad\forall i
 \end{aligned}\right.
\]

Soit de $\overline{\mathrm m}$
\[
 \left|\begin{aligned}
 g&=\sum_J\beta_j B_j&\beta_j>0,\quad B_j\in S_0\quad (B_j)\ \text{tot ord.}\\
 t&\in G,&B_j=\varepsilon_t*\check B_j
 \end{aligned}\right.
\]
On a
\[
 f*g=\sum_{(i,j)\in I\times J}\alpha_i\beta_j\, A_i*B_j
\]

\deleted{Supposons que $\exists\ (\unsure{C_k|C_h|G_n})$}

On a $A_i*B_j\in S$, d'où
\[
 E_*(A_i*B_j)\subset S_0
\]

\uuline{Supposons} $\bigcup E_{\alpha}(A_i*B_j)$ \uline{tot} ordonné

ou seul\textsuperscript{t}
\[
 \bigcap_{C\in\bigcup_{ij}E_\alpha(A_i*B_j)} C\ne\varnothing
\]

\diagram{\begin{tikzcd}[ampersand replacement=\&,column sep=small]
 0\arrow[r] \& {}_2G\arrow[r] \& G\arrow[r,"2"] \& G
\end{tikzcd}}
% archive sheet number: 21
\scanpage{24}
\uline{Jeux de position}

a) $C$ ensemble des ``\uline{configurations}'' ou ``\uline{positions}''

b) $\mathcal R\subset C\times C$, interprété comme $x\mapsto\mathcal R(x)$ \quad $C\longrightarrow\mathcal P(C)$

d'où
\[
 C_0=\{x\in C\mid\mathcal R(x)=\varnothing\}\subset C\qquad(\text{conf. \uline{terminales}})
\]

c) $J$ ensemble des ``\uline{joueurs}''

d) $C\setminus C_0\xrightarrow{\alpha}J$ \quad On dit que (si $x\in C\setminus C_0$) $j=\alpha(x)$ est le joueur qui a \uline{l'initiative} dans la position $x$.
\margin{$C(j)=\{x\in C\setminus C_0\mid\alpha(x)=j\}$\\$\bar C(j)=(C\setminus C_0)\setminus C(j)$}

e) $J\longrightarrow\mathcal P(C_0)$ \qquad $j\longrightarrow G_0(j)$ (ens. des \uline{positions terminales gagnantes pour} $j$)

\uline{Partie} $(x_0,x_1,\ldots,x_n,\ldots)$ suite finie ou infinie de positions, telles que pour deux termes consécutifs $x_i,x_{i+1}$, on ait
\[
 x_{i+1}\in\mathcal R(x_i)
\]
% slanted
\margin{ne dépend que de $C,\mathcal R$}

\uline{Prolongement} des parties -- parties non prolongeables ou ``\uline{terminées}'' : ce sont les parties qui sont soit infinies, soit finies, avec \deleted{\unsure{comme}}\inserted{la} position finale $x_n\in C_0$.

\margin{Composition des parties (Associatif etc -- on a une ``catégorie des positions'' $\ldots$ )}

\uline{Parti\deleted{\unsure{es}}e gagnée pour $j$} : partie finie (non prolongeable) dont la position \deleted{fi}nale $x_n$ est $\in G_0(j)$.
% slanted
\margin{ne dépend que de $C,\mathcal R$ et $G_0(j)\subset C_0$}

\margin{\uline{NB} \deleted{Si} $XY$ est gagnée \unsure{par|pour} $j$, \deleted{\unsure{alors}}$\Longleftrightarrow Y$ est gagnée \unsure{par|pour} $j$.}

\uline{Stratégie \inserted{$\Sigma$} pour $j$} \inserted{(ou $j$-stratégie) :} Application \inserted{$\Sigma:$} $C(j)\longrightarrow\mathcal P(C)$ telle que $\forall x\in C(j)$, on ait
% archive sheet number: 23
\scanpage{25}
\[
 \Sigma(x)\subset R(x),\qquad\Sigma(x)\ne\varnothing .
\]

Cette notion ne dépend que de $C,\mathcal R,\ C(j)\subset C\setminus C_0$

\uline{NB} On pose $\Sigma(x)=R(x)$ si $x\in\bar C(j)$. Donc on peut

\deleted{Stratégie pour $j$ \unsure{gagnante pour} $x_0\in C$ :}

\inserted{interpréter $\Sigma$ comme une application $C\setminus C_0\longrightarrow\mathcal P^*(C)$, avec $\Sigma(x)\subset R(x)$ et $\Sigma(x)\ne\varnothing\quad\forall x\in C\setminus C_0$, et $\Sigma(x)=R(x)$ si $x\in\bar C(j)$.}

Partie jouée en accord avec une stratégie $\Sigma$ (ou \uline{compatible} avec $\Sigma$, ou $\Sigma$-partie) : cette notion est stable par composition. (\unsure{Définit|On a une} sous-catégorie ayant $\overline{\mathrm m}$s objets, mais a priori moins de flèches)

\inserted{\uline{Lemme} Soit \deleted{\unsure{$\Sigma$ gagnante}} $J'\subset J$, $\forall j\in J'$ soit $\Sigma_j$ une $j$-stratégie. Soit $x_0\in C$. Alors $\exists$ partie maximale d'origine $x_0$, comp. avec les $\Sigma_j$ $(j\in J)$.}

Stratégie $\Sigma$ (pour $j$) \uline{gagnante} \uline{pour} $x_0\in C$ (ou \uline{position} $x_0$ \uline{gagnante} pour une $j$-stratégie $\Sigma$) : toute partie maximale commençant en $x_0$, et jouée suivant $\Sigma$, est gagnée par $j$ (i.e. est finie, et son extrémité $x_n\in C_0$ est $\in G_0(j)$).
% slanted
\margin{ne dépend que de $C,\mathcal R,\unsure{C(j)|\bar C(j)},G_0(j)$}

\uline{Position \inserted{$x_0\in C$} gagnante pour} $j$ : telle qu'il existe une $j$-stratégie qui soit gagnante pour $x_0$. Soit
\[
 G(j)
\]
l'ens des \unsure{ces|ses} positions. On a.
\[
 G(j)\cap C_0=G_0(j)
\]

[ \uline{Proposition} Soit $J'\subset J$. \deleted{\unsure{Soit}} Si $\displaystyle\bigcap_{j\in J'}G(j)\ne\varnothing$, alors $\displaystyle\bigcap_{j\in J}G_0(j)\ne\varnothing$. En d'autres termes, si $\displaystyle\bigcap_{j\in J'}G_0(j)=\varnothing$, alors $\displaystyle\bigcap_{j\in J'}G(j)=\varnothing$ i.e. une $x_0\in C$ ne peut être gagnante pour tous les $j\in J'$.

Immédiat par le lemme. ]
% archive sheet number: 24
\scanpage{26}
\uline{Proposition} Soient $\Sigma$ une $j$-stratégie, \deleted{Supposons}\inserted{$x_0\in C$} gagnante pour $\Sigma$, $X=(x_0,x_1,\text{--},x_n\ldots)$ une partie (d'origine $x_0$) jouée suivant $\Sigma$, alors toutes les positions $x_i$ sont gagnantes pour $\Sigma$.

Soit en effet $Y=(y_0=x_i,y_1,\text{--},y_m,\text{--})$ une partie \inserted{maximale} jouée suivant $\Sigma$, \deleted{\unsure{on}} alors $XY$ est jouée suivant $\Sigma$ et est d'origine $x_0$ gagnante pour $\Sigma$, donc est \deleted{\unsure{est|partie} finie \unsure{Cdà|telle} $Y$ est finie}\inserted{gagnée par} donc $Y$ est gagnée par $j$.

\uline{Th} On a \deleted{\unsure{Supposé}}
\[
 \begin{aligned}
 \text{a)}\quad G(j)\cap C(j)&=\{x_0\in C(j)\mid R(x_0)\cap G(j)\ne\varnothing\}\\
 \text{b)}\quad G(j)\cap\bar C(j)&=\{x_0\in\bar C(j)\mid R(x_0)\subset G(j)\}
 \end{aligned}
\]
[ et pour mémoire
\[
 G(j)\cap C_0=\unsure{C_0|G_0}(j)\qquad ].
\]

\uline{Dém} \inserted{a) On suppose $x_0\in C(j)$.}

\deleted{\unsure{a)}} 1) $\subset$ Soit $x_0\in G(j)\deleted{\unsure{\cap C(j)}}$, \inserted{soit} $\Sigma$ une stratégie pour $j$ gagnante pour $x_0$, considérons $x_1\in\Sigma(x_0)$, alors $(x_0,x_1)$ est une $\Sigma$-partie, donc par la prop. précédente $x_1$ est gagnante pour $\Sigma$, donc gagnante, donc $x_1\in R(x)\cap G(j)$, donc $R(x)\cap G(j)\ne\varnothing$.

2) $\supset$ Soit $x_1\in R(x)\cap G(j)$. Prouvons que $x_0\in G(j)$. \deleted{On \unsure{défi}} Soit $\Sigma_1$ une $j$-stratégie $x_1$-gagnante. Si $x_0$ est gagnante pour $\Sigma_1$, on a $x_0\in G(j)$, OK. Sinon, définissons une $j$-stratégie $\Sigma_0$ $x_0$-gagnante ainsi
\[
 \begin{aligned}
 \Sigma_0(x_0)&=\{x_1\}\\
 \Sigma_0(x)&=\Sigma_1(x)\quad\text{si}\quad x\ne x_0
 \end{aligned}
\]
% archive sheet number: 25
\scanpage{27}
C'est une $j$-stratégie, \deleted{prouvons}\inserted{Prouvons} qu'elle est gagnante pour $x_0$\deleted{ou}. Une partie maximale jouée suivant $\Sigma_0$ \unsure{commence|commençant} avec $(x_0,x_1)$, soit $(x_0,x_1,\text{--},x_n,\text{--})$, \deleted{les $x_i$ $(i\geq1)$ sont gagnants pour $\Sigma_1$} prouvons qu'elle est gagnée par $j$. Prouvons \unsure{que tous} les $x_i$ $(i\geq1)$ sont $\ne x_0$. Sinon, \deleted{\unsure{on aurait|si $i$ est le}}\inserted{si \deleted{\unsure{si}} $i$ est le} plus petit indice $i\geq1$ tel que $x_i=x_0$, \deleted{pour un $i\geq1$,} alors \deleted{par la prop. \inserted{\unsure{appliquée|impliquerait}} à} $(x_1,\ldots,x_i=x_0)$ \inserted{serait compatible avec $\Sigma_1$,} on aurait que $x_{i}\inserted{{}=x_0}$ est gagnante pour $\Sigma_1$, contrairement à la supposition que $x_0$ non gagnante pour $\Sigma_1$.

Comme les $x_i$ $(i\geq1)$ sont $\ne x_0$, la partie $(x_1,x_2,\deleted{\unsure{x_0}}\text{--})$ est une $\Sigma_1$-partie, donc elle est gagnée par $j$, donc $(x_0,x_1,\text{--})$ aussi, cqfd.

b) On suppose $x_0\in\overline{C(j)}$.

1) $\subset$. Supposons $x_0\in G(j)$, prouvons $R(x)\subset G(j)$. Soit $x_1\in\mathcal R(x)$, prouvons $x_1\in G(j)$. En effet, \inserted{si $\Sigma_0$ est une $j$-stratégie $x_0$-\unsure{gagn.},} $(x_0,x_1)$ est une \inserted{$\Sigma_0$-}partie, donc par la prop, $x_1$ est gagnante pour $\Sigma_0$, donc $x_1\in G(j)$.

2) $\supset$. Supposons $R(x_0)\subset G(j)$, prouvons $x_0\in G(j)$.

$\forall x_1\in R(x_0)$, soit $\Sigma_{x_1}$ une $j$-stratégie $x_1$-gagnante. Choisissons un \uline{bon ordre} sur $\deleted{E}R(x_0)$. On va définir une $j$-stratégie $\Sigma_0$ $x_0$-gagnante ainsi.

$\forall x\in C$, soit
\[
 E(x)=\{x_1\in R(x_0)\mid\exists\ \Sigma_{x_1}\text{-partie de }x_{\unsure{1|0}}\text{ à }x\}
\]
% archive sheet number: 26
\scanpage{28}
Soit
\[
 C'=\{x\in C\mid E(x)\ne\varnothing\}
\]
et pour $x\in C'$, posons
\[
 \varepsilon(x)=\text{plus petit élément de } E(x)
\]
(donc $\varepsilon:C'\longrightarrow R(x_0)$). Posons
\[
 \Sigma_0(x)=\begin{cases}
 R(x) & \text{si } \deleted{\unsure{x}}x\in\bar C(j)\cup\complement C'\\
 \Sigma_{\varepsilon(x)}(x)&\text{si }x\in C(j)\cap C'
 \end{cases}
\]
Donc $\Sigma_0$ est une $j$-stratégie. On va prouver qu'elle est gagnante pour $x_0$.

\uline{Lemme} \deleted{\unsure{On a|On à}} \inserted{Si $x\in C'$, $x$ est gagnant pour $\Sigma_{\varepsilon(x)}$ (donc gagnant pour $j$ -- donc $C'\subset G(j)$) et \deleted{\unsure{si}}} \deleted{$C'\subset G(j)$, et si $x\in C'$, $y\in\Sigma_0(x)$,} alors \deleted{$\Sigma_0(x)\subset C'$ et} pour tt $y\in\Sigma_0(x)$, on a $E(y)\ni \deleted{\unsure{E}}\varepsilon(x)$ (donc $E(y)\ne\varnothing$ donc $y\in C'$, donc $\Sigma_0(x)\subset C'\inserted{\subset G(j)}$) et par suite
\[
 \varepsilon(y)\leq\varepsilon(x).
\]

\uline{\unsure{Dém.|Dém}} \deleted{\unsure{$C'$|1\textsuperscript{o}}} Soit $x\in C'$ donc $E(x)\ne\varnothing$, alors si $x_1\in E(x)$ $\exists$ \deleted{partie} $\Sigma_{x_1}$-partie de $x_1$ à $x$, donc par la prop. (comme $\Sigma_{x_1}$ est $x_1$-\unsure{gagnante|gagnant}), $x$ est $\Sigma_{x_1}$-gagnante, donc à fortiori $x\in G(j)$.

Soit $y\in\Sigma_0(x)=\Sigma_{\varepsilon(x)}(x)$
\[
 (\underset{\substack{\Vert\\ \varepsilon(x)}}{x_1},x_2,\ldots,x_n=x,x_{\unsure{n+1|n-1}}=y)
\]
est une $\Sigma_{\varepsilon(x)}$-partie, donc \deleted{\unsure{$E$|$E\in$}} $\varepsilon(x)\in E(y)$, OK.

Prouvons que $\Sigma_0$ est gagnante pour $x_0$. Soit une partie maximale commençant par $x_0$
% archive sheet number: 27
\scanpage{29}
et comp. avec $\Sigma_0$
\[
 (x_0,x_1,x_2,\ldots)
\]
Par le lemme, on voit de proche en proche \deleted{Cette partie est} que \deleted{$x_1\in G(j)$, $x$} \inserted{(puisque $x_1\in C'$) \unsure{que,|que}} les $x_i$ $(i\geq1)$ sont $\in C'\deleted{\unsure{\subset G(j)|\in G(j)}}$ donc $j$-gagnantes. Pour montrer que la partie est gagnée par $j$, il suffit de prouver qu'elle est finie. \deleted{Or} \inserted{Supposons l'inverse.} Les $\varepsilon(x_i)$ $(i\geq1)$ forment une suite décroissante d'éléments de $\deleted{\unsure{P|R}}R(x_0)$, donc stationnaire. Soit $n$ tel que $\varepsilon(x_m)=\varepsilon(x_n)\inserted{(=\xi)}$ pour $m\geq n$. \deleted{\unsure{Par}} Alors par définition de $\Sigma_0$, la partie $(x_n,x_{n+1},\ldots,\,)$ est une partie comp. avec $\Sigma_\xi$. Comme $x_n$ est gagnant pour $\Sigma_{\unsure{\xi|j}}$ (par le lemme), cette partie est finie, absurde.

Stratégie \uline{contre $j$} (\deleted{stratégie} \unsure{relation|relative|relative à} : $\bar C(j)$), i.e. $\Sigma'$ telle que $\Sigma'(x)=R(x)$ si $x\in C(j)$ (et $\Sigma'(x)\subset R(x)$, $\Sigma'(x)\ne\varnothing$ si $x\in\bar C(j)$) \unsure{admissible} pour $x_0$ (tte partie maximale d'origine $x_0$ jouée suivant $x_0$ est \uline{non gagnée} par $j$, i.e. est ou bien infinie, ou bien a une position terminale $\notin \unsure{G_0(j)|\bar G_0(j)|C_0(j)}$).

\uline{\deleted{Corollaire}} \inserted{Théorème} Soit $\Sigma'$ la stratégie contre $j$ définie par
% archive sheet number: 28
\scanpage{30}
\[
 \Sigma'(x)=\begin{cases}
 R(x) \text{ si } & x\in C(j)\cup G(j)\\
 R(x)\cap\complement G(j) \text{ si } & x\in\bar C(j)\cap\complement G(j)
 \end{cases}
\]
(C'est une stratégie contre $j$, par $[\,\mathrm b)\ \supset\,]$. \deleted{\unsure{Al}}Alors Pour $x\in C$, les conditions suivantes, sont équivalentes

a) $x_0\notin G(j)$

b) $x_0$ est admissible contre $j$

c) $x_0$ est admissible pour $\Sigma'$ contre $j$.

\uline{Dém} $\mathrm c)\Longrightarrow\mathrm b)$ trivial, $\mathrm b)\Longrightarrow\mathrm a)$ par le lemme, Prouvons $\mathrm a)\Longrightarrow\mathrm c)$. Si on a une partie $(x_0,x_1,\ldots)$ d'\deleted{\unsure{ori}}origine $x_0$ jouée en accord avec $\Sigma'$, alors par définition de $\Sigma$, on voit de proche en proche que $x_i\in\complement G(j)\quad\forall i$, donc la partie ne peut être gagnée par $j$ !

\uline{Corollaire} Supposons $J=\{j,j'\}$, $j\ne j'$, et \margin{$\vert$} $G_0(j)\cap G_0(j')=\varnothing$. Soit \deleted{$\unsure{N}$}
\[
 N=\{x\in C\mid x\text{ admissible contre }j,\text{ et contre }j'\}.
\]
\margin{$\big\vert$} Alors $C$ est réunion disjointe de $G(j)$, $G(j')$ et $N$.
% archive sheet number: 29
\scanpage{31}
\deleted{\unsure{Suppose|Suppos|Supposons}} \uline{NB} $G(j),G(j')$ et $N(j)$ sont connus quand on connaît (en plus des \inserted{parties} $G_0(j)$ et $G_0(j')$ de $C_0$) les ensembles
\[
 \left\{\begin{aligned}
 G^*&=\{x\in C\setminus C_0\mid x\in G(\alpha(x))\}\\
 P^*&=\{x\in C\setminus C_0\mid x\in G(\bar\alpha(x))\}
 \end{aligned}\right.
\]
où $\bar\alpha(x)$ désigne l'unique $j'\in J$ tel que $j'\ne j=\alpha(x)$.

On dit qu'une position $x\in C\setminus C_0$ qui est $\in G^*$ \inserted{(resp. $\in P^*$)}, i.e. \deleted{\unsure{telle|tels|telles}} \inserted{gagnante} \deleted{pour} le joueur qui a l'initiative (resp. pour le joueur qui n'a pas l'initiative.) est \uline{gagnante} (resp \uline{perdante}) ; on dit que c'est une \uline{position nulle} \deleted{autrement} si elle n'est ni gagnante, ni perdante, on pose
\[
 N^*=(C\setminus C_0)\setminus(G^*\cup P^*)=N\cap(C\setminus C_0)
\]
(où $N$ est défini ds le corollaire ci-dessus).

On aura
\[
 \left\{\begin{aligned}
 G(j)\cap C(j)&=G^*\cap C(j)\\
 G(j)\cap\bar C(j)&=G(j)\cap C(j')=P^*\cap C(j')\\
 G(j)\cap C_0&=G_0(j)
 \end{aligned}\right.
\]
ce qui montre que \inserted{les} $G(j)$ $(j\in G)$ se réconstituent par la connaissance de $G^*,P^*$ et des $G_0(j)$ $(j\in J)$.

Supposons qu'on ait une application
\[
 \hat\alpha:C\longrightarrow J
\]
prolongeant $\alpha$, et qu'on ait, pour \unsure{tout|tt} $x\in C$ et $y\in R(x)$
\[
 \hat\alpha(y)\ne\hat\alpha(x)
\]
% archive sheet number: 30
\scanpage{32}
On dit alors que le jeu \deleted{\unsure{est|de}} de position à deux joueurs est un jeu \uline{alterné} (les deux joueurs jouent alternativement). Pour qu'une application $\hat\alpha$ existe prolongeant $\alpha$ et ayant les propriétés précédente, il f. et s. que les deux conditions suivantes soient satisfaites

1) $\forall x,y\in C$, tels que $y\in R(x)$, $y\notin \deleted{\unsure{C\setminus|C_0}} C_0$, on a $\alpha(y)\ne\alpha(x)$

2) $\forall y\in C_0$, soit $T(y)\deleted{\unsure{\in|x}}=\{x\in C\mid y\in R(x)\}\subset\deleted{\unsure{C_0}}C\setminus C_0$, alors $\alpha\mid T(y)$ est constante.

Sous les conditions de 2), on aura alors nécessairement, si $T(y)\ne\varnothing$ et $\alpha(T(y))=\{j\}$, $\hat\alpha(y)=j'\deleted{\unsure{\ne j|j}}$, où $j'$ est l'élément de $J$ qui est $\ne j$. Cela montre que $\hat\alpha$ est bien déterminé sur l'ens \deleted{\unsure{de}} des $y\in C_0$ tels que $T(y)\ne\varnothing$, donc aussi sur l'ens des positions \uline{non} \uline{isolées} (une position \inserted{$y$} est dite \uline{isolée} si elle n'a pas de position antécédente ni suivante, i.e. s'il n'existe pas de $x\in C$ tel qu'on ait $y\in R(x)$ ou $x\in R(y)$). On peut attribuer \unsure{à|:} $\hat\alpha$ des valeurs arbitraires sur les positions isolées.

Quand on a défini \inserted{un} $\hat\alpha$, on peut définir
\[
 \begin{aligned}
 G&=\{x\in C\mid x\in G(\alpha(x))\}\\
 P&=\{x\in C\mid x\in G(\overline{\hat\alpha}(x))\}
 \end{aligned}
\]
et on \deleted{$\unsure{N=\complement(G\cup P)|N=C\setminus(G\cup P)}$ \unsure{aura encore|leur réunion}} aura
% archive sheet number: 31
